\section{The diagonal certificate and the proximal stage}\label{app:proximal}

This appendix proves Lemma~\ref{lem:diagcert}, Lemma~\ref{lem:proximal},
Theorem~\ref{thm:diagdiscovery} and Proposition~\ref{prop:sshard}. As in
Section~\ref{sec:diagcert}, $X$ is a continuous box and $F$ a rational
quadratic.

\begin{proof}[Proof of Lemma~\ref{lem:diagcert}]
(i) Both sides are quadratic polynomials in $x$ that vanish at $\bar x$.
Their Hessians are $H$ and $H+\Lambda-\Lambda=H$. The $i$th partial
derivative of the right side at $\bar x$ is
$\frac12\lambda_i(\bar x)(u_i+\ell_i-2\bar x_i)$. This equals
$|\partial_iF(\bar x)|=\partial_iF(\bar x)$ when $\bar x_i=\ell_i$, equals
$-|\partial_iF(\bar x)|=\partial_iF(\bar x)$ when $\bar x_i=u_i$, and is
$0=\partial_iF(\bar x)$ at an interior coordinate. Two quadratics with equal
value, gradient and Hessian at one point coincide.

(ii) On $X$ both terms on the right of \eqref{eq:diagid} are nonnegative,
so $\bar x$ is optimal, and $x\in X$ is optimal exactly when both vanish. A
positive semidefinite quadratic form vanishes exactly on the kernel of its
matrix.

(iii) For $x\in X$ the bracket in $\Psi$ is at most $F(x)$, so
$\Psi(\lambda)\le\OPT$. If (a) holds at $\bar x$, then by (i) the bracket
with $\lambda=\lambda(\bar x)$ equals
$F(\bar x)+\frac12(x-\bar x)^{\T}(H+\Lambda(\bar x))(x-\bar x)$, whose
infimum over $\R^n$ is $F(\bar x)=\OPT$; this gives (c). Suppose (c) holds
for some $\lambda$, and let $\bar x\in\mathcal S$ be arbitrary. Write
$\mathcal L(x)$ for the bracket. Then
$\OPT=\Psi(\lambda)\le\mathcal L(\bar x)\le F(\bar x)=\OPT$. Hence
$\lambda_i(\bar x_i-\ell_i)(u_i-\bar x_i)=0$ for all $i$, and $\bar x$
minimizes the quadratic $\mathcal L$ over $\R^n$. So
$\nabla^2\mathcal L=H+\diag(\lambda)\succeq0$ and
$0=\partial_i\mathcal L(\bar x)=\partial_iF(\bar x)+\lambda_i\bigl(\bar x_i-\tfrac{\ell_i+u_i}2\bigr)$.
At an interior coordinate, complementarity gives
$\lambda_i=0=\lambda_i(\bar x)$. At $\bar x_i=\ell_i$ the stationarity
equation gives $\lambda_i=2\partial_iF(\bar x)/(u_i-\ell_i)=\lambda_i(\bar x)$,
and similarly at $\bar x_i=u_i$. Thus $\lambda=\lambda(\bar x)$ and
$H+\Lambda(\bar x)\succeq0$. Since $\bar x\in\mathcal S$ was arbitrary, this
proves (b) and the uniqueness and invariance statements. Finally (b) implies
(a) because $\mathcal S\ne\emptyset$.
\end{proof}

\begin{proof}[Proof of Lemma~\ref{lem:proximal}]
\emph{Constants.} Since $\mu\ge2$, $\theta\le1/4$, and
$4^{-\mu}\hat\kappa\le1/8$ is $8\hat\kappa\theta^2\le1$. If $\mu=2$, then
$\theta^{-1}=4\le6\sqrt{\hat\kappa}$. Otherwise minimality of $\mu$ gives
$4^{1-\mu}\hat\kappa>1/8$, so $\theta^{-1}=2^\mu<2\sqrt{8\hat\kappa}
<6\sqrt{\hat\kappa}$.

(i) Each endpoint of $X'_i$ is within $\varrho h$ of $c_i$, so by
Lemma~\ref{lem:graded}(b) each side of $G_i$ has at most
$\lceil(4/\theta)\ln(1+\theta\varrho)\rceil$ nodes besides $c_i$. Now
$\theta\varrho\le2\theta\sqrt{\hat\kappa n}+2\theta\le\sqrt{n/2}+\frac12$,
because $\theta\sqrt{\hat\kappa}\le1/\sqrt8$ and $\theta\le\frac14$, and
$\frac32+\sqrt{n/2}\le n+2$. Hence
$\abs{G_i}\le1+2\lceil(4/\theta)\ln(n+2)\rceil\le3+(8/\theta)\ln(n+2)$. With
$3\le3/(4\theta)$, $8\ln(n+2)\le5.6\log_2(n+2)$ and
$\lceil\log_2(n+2)\rceil\ge2$, this is at most
$7\theta^{-1}\lceil\log_2(n+2)\rceil\le K_\theta$.

(ii) Let $s^*\in\mathcal S$ be nearest to $c$ (it exists since $\mathcal S$
is compact), and put $\delta=\norm{s^*-c}$, so $\delta^2\le4\hat\kappa nh^2$. Each
$|s^*_i-c_i|\le\delta\le2\sqrt{\hat\kappa n}\,h\le\varrho h$, so $s^*\in X'$.
For each $i$ choose a grid interval $J_i$ of $G_i$ that contains $s_i^*$ and
whose endpoint $a_i$ nearer to $c_i$ satisfies $|a_i-c_i|\le|s_i^*-c_i|$:
if $s_i^*>c_i$, the interval $[a_i,a_i']$ with $a_i<s_i^*\le a_i'$; if
$s_i^*<c_i$, symmetrically; if $s_i^*=c_i$, any interval with endpoint
$c_i$. By Lemma~\ref{lem:graded}(a) at the node $a_i$,
$w(J_i)\le h+\theta|s_i^*-c_i|$. Let $C=\prod_iJ_i$ and
$\tilde F(y)=F(y)+\eta\norm{y-c}^2$, which has the upper coordinate
curvatures $L_i+2\eta$. The proof of Proposition~\ref{prop:cellwise}(a) uses
only Definition~\ref{def:curvature}, so it applies to $\tilde F$ and gives a
vertex $v$ of $C$ with
$\tilde F(v)-\sum_i\frac{L_i+2\eta}8w(J_i)^2\le\tilde F(s^*)=\OPT+\eta\delta^2$.
Since $v_i$ is an endpoint of $J_i$, $d_i(v_i)\ge L_iw(J_i)^2/8$, and so
\[
Q_\eta(y^+)\le Q_\eta(v)\le\tilde F(v)-\sum_i\frac{L_i}8w(J_i)^2
\le\OPT+\eta\delta^2+\frac\eta4\sum_iw(J_i)^2 .
\]
With $w(J_i)^2\le2h^2+2\theta^2(s_i^*-c_i)^2$ this gives
\[
Q_\eta(y^+)\le\OPT+\eta\Bigl[\bigl(1+\tfrac{\theta^2}2\bigr)\delta^2
+\tfrac{nh^2}2\Bigr].
\]
By Lemma~\ref{lem:graded}(a), $w_i(v)\le h+\theta|v-c_i|$ for every node $v$,
so $d_i(v)\le\frac L4h^2+\eta|v-c_i|^2$ and
$D(y^+)\le\frac L4nh^2+\eta\norm{y^+-c}^2$. Therefore
\begin{align*}
F(y^+)-\OPT&=Q_\eta(y^+)-\OPT+D(y^+)-\eta\norm{y^+-c}^2\\
&\le\frac{Lnh^2}4\Bigl[1+4\hat\kappa\theta^2\bigl(1+\tfrac{\theta^2}2\bigr)
+\tfrac{\theta^2}2\Bigr]\le\frac{Lnh^2}4\cdot\frac{99}{64},
\end{align*}
using $\delta^2\le4\hat\kappa nh^2$, $\eta=L\theta^2/4$,
$4\hat\kappa\theta^2\le\frac12$ and $\theta^2\le\frac1{16}$. Set growth gives
$\dist(y^+,\mathcal S)^2\le(F(y^+)-\OPT)/g_S\le\frac{99}{256}(L/g_S)nh^2\le
\frac{99}{256}\hat\kappa nh^2$, since $L/g_S\le\kappa_S\le\hat\kappa$.

(iii) By induction on $j$, the center of stage $j$ satisfies
$\dist(c^{(j)},\mathcal S)^2\le4\hat\kappa nh_j^2$, so (ii) applies at stage
$j$. For $j=0$, $\dist(\ell,\mathcal S)^2\le ns^2=nh_0^2\le4\hat\kappa nh_0^2$.
If (ii) applies at stage $j$, then
$\dist(c^{(j+1)},\mathcal S)^2=\dist(y^{(j)},\mathcal S)^2
\le\frac{99}{256}\hat\kappa nh_j^2<2\hat\kappa nh_{j+1}^2$.

(iv) Let $\Delta$ be the common denominator of Section~\ref{sec:exact-data};
it also clears $s$ and $L$. By induction on $j$, every node of stage $j$ lies
in $\omega_j^{-1}\Z$ with $\omega_j=\Delta2^{j+\mu K_\theta}$: the center is
$\ell$ or a node of stage $j-1$; the box endpoints are $\ell_i$, $u_i$ or
$c_i\pm\varrho h_j$ with $h_j=s2^{-j}$; an unclipped node is $c_i\pm t_k$ with
$t_k=h_j\bigl[(2^\mu+1)^k-2^{\mu k}\bigr]/2^{\mu(k-1)}$ and $k<K_\theta$; and a
clipped node is a box endpoint. Hence, at every grid point, $\Delta\omega_j^2F$
is an integer (Section~\ref{sec:exact-data}), and so are $8\Delta\omega_j^2d_i$
and $4^{\mu+1}\Delta\omega_j^2\eta\norm{y-c}^2$. All table entries are
therefore integers over the common denominator $8\cdot4^\mu\Delta\omega_j^2$,
and their magnitudes are at most $2^{O(I)}$ times this denominator because
all points lie in $X$. Messages are sums of at most $\abs{\mathcal A}+2n$
such values, so all numbers have bit length $O(I+j+\mu K_\theta)$. The count
of table operations is that of Lemma~\ref{lem:dp} with at most $K_\theta$
nodes per coordinate.
\end{proof}

\begin{proof}[Proof of Theorem~\ref{thm:diagdiscovery}]
(a) DISC (Algorithm~\ref{alg:disc}) stops only when the hypotheses of
Lemma~\ref{lem:diagcert}(ii) hold, and it checks them in exact rational arithmetic: the KKT conditions
are sign tests, and positive semidefiniteness of $H+\Lambda(\bar x)$ is
decided by symmetric Gaussian elimination that rejects a negative pivot and
a zero pivot with a nonzero entry in its row. By
Lemma~\ref{lem:diagcert}(ii), $\bar x\in\mathcal S$, $\OPT=F(\bar x)$
and~\eqref{eq:diagset} holds. So an accepted $\bar x$ is a point of
$\mathcal S$ with $H+\Lambda(\bar x)\succeq0$, and then $F\in\mathfrak D$
by Lemma~\ref{lem:diagcert}(iii). Hence DISC does not stop if
$F\notin\mathfrak D$.

(b) Remark~\ref{rem:setgrowth} gives some $g_S>0$. Let $\hat\kappa\ge\kappa_S$
be a guess. By Lemma~\ref{lem:proximal}(iii) with $h_{J_{\hat\kappa}}^2
=s^24^{-J_{\hat\kappa}}\le\tau^2/(4\hat\kappa n)$,
$\dist(y,\mathcal S)^2\le\frac{99}{256}\cdot\frac{\tau^2}4<\frac{\tau^2}4$. By
Lemma~\ref{lem:snap}, REC$(y)$ (Algorithm~\ref{alg:rec}) does not fail, and every point it can return,
in particular $\bar x$, lies in $\mathcal S$; so $\bar x$ is a box KKT point.
Since $F\in\mathfrak D$, Lemma~\ref{lem:diagcert}(iii) gives
$H+\Lambda(\bar x)\succeq0$ for every $\bar x\in\mathcal S$, and the test
accepts, whichever optimal component contains $\bar x$. The guesses are
powers of two, so the first guess $\hat\kappa\ge\kappa_S\ge1$ satisfies
$\hat\kappa\le2\kappa_S$.

(c) By (b), at most $2+\log_2\kappa_S$ guesses run, each with
$\hat\kappa\le2\kappa_S$. For each,
$J_{\hat\kappa}+1=O(\log(\hat\kappa ns^2/\tau^2))+1=\poly(I)+O(\log\hat\kappa)$,
because $\tau=1/(4nR)$ and $\log R$ is polynomial in $I$. By
Lemma~\ref{lem:proximal}(i) and~(iv), a stage costs
$O(p(\abs{\mathcal A}+n+N)K_\theta^p)$ operations on numbers of bit length
$O(I+J_{\hat\kappa}+\mu K_\theta)$, where $\mu=O(\log(2\hat\kappa))$ and
$K_\theta\le48\sqrt{\hat\kappa}\lceil\log_2(n+2)\rceil$. By~\eqref{eq:logabsorb}
with $n$ in place of $n_P$, $K_\theta^p\le2(48p\sqrt{\hat\kappa})^p(n+2)$, and
$\mu K_\theta\le c\sqrt{\hat\kappa}\log(2\hat\kappa)\,I$ for an absolute
constant $c$. Every factor that depends only on $p$ and $\hat\kappa$ is
nondecreasing in $\hat\kappa$; its value at $\hat\kappa=2\kappa_S$ enters
$f'$, and for fixed $p$ these values are nondecreasing and polynomial in
$\kappa_S$. The remaining factors are polynomial in $I$ with an absolute
exponent. REC uses the original data, so its linear system has size
polynomial in $I$, and exact linear programming returns a vertex of
polynomial bit length in polynomial time
\cite[Theorem~6.4.12]{GrotschelLovaszSchrijver1988}. The KKT and
semidefiniteness tests are polynomial.
\end{proof}

\begin{proof}[Proof of Proposition~\ref{prop:sshard}]
$F\ge0$ on the box, and $F=0$ exactly when all squares and all products
$x_i(1-x_i)$ vanish, that is, for binary $x$ with partial sums $\xi$ and
$\xi_m=a_0$; the partial sums lie in $[-A,A]$. If no solution exists, $F>0$
on the compact box. In the feasible case let $\hat z=(\hat x,\hat\xi)$ be
optimal. The sum of squares is a quadratic $\Phi$ with $\Phi(\hat z)=0$ and
$\nabla\Phi(\hat z)=0$, because each square is of an affine function
vanishing at $\hat z$; so $\Phi(z)=\frac12(z-\hat z)^{\T}H_{\mathrm{sq}}
(z-\hat z)$ with $H_{\mathrm{sq}}=\nabla^2\Phi\succeq0$. Also
$\sum_ix_i(1-x_i)=\frac12\sum_i2(x_i-0)(1-x_i)$. So
$F(z)-F(\hat z)=\frac12(z-\hat z)^{\T}H_{\mathrm{sq}}(z-\hat z)
+\frac12\sum_i2x_i(1-x_i)$. At $\hat z$ the gradient is $1-2\hat x_i=\pm1$ in
$x_i$ and $0$ in $\xi$, so $\hat z$ is a box KKT point with
$\lambda_i(\hat z)=2$ on $x$ and $0$ on $\xi$, and
$H+\Lambda(\hat z)=H_{\mathrm{sq}}\succeq0$. Thus $F\in\mathfrak D$.

Subset Sum with binary-encoded integers is NP-complete
\cite{GareyJohnson1979}, also when restricted to targets with
$\abs{a_0}\le A$, because other targets have no solution; and the instance $F$
has size polynomial in that of the Subset Sum instance. (i) An algorithm as in (i), stopped at its
polynomial time bound and followed by a check that its output $z$ is
feasible with $F(z)=0$, would decide Subset Sum. (ii) If such a polynomial
$\pi$ existed, then by Theorem~\ref{thm:diagdiscovery}(c), whose bound is
nondecreasing and polynomial in $\kappa_S$ for fixed $p=3$, DISC would stop
within a polynomial time bound on every feasible instance; stopping DISC at that bound and
checking whether it has output $\OPT=0$ would decide Subset Sum.
\end{proof}

\begin{remark}[A false guess]\label{rem:falseguess}
If $\hat\kappa<\kappa_S$, PROX$(\hat\kappa)$ (Algorithm~\ref{alg:prox}) may
return the same nonoptimal
point at every stage. For $F=x^2+y^2-3xy+\frac{63}{128}(x+y)$ on $[0,1]^2$, put
$t=x+y$; then $F\ge t^2-\frac54t^2+\frac{63}{128}t$, a concave function of
$t\in[0,2]$, so $(1,1)$ is the unique optimizer with $\OPT=-\frac1{64}$. The
origin is a box KKT point with $F=0$, so every set-growth constant
satisfies $g_S\le(F(0)-\OPT)/\dist(0,\mathcal S)^2=\frac1{128}$, and
$\kappa_S\ge256$ since $L=2$. Exact computation shows that
PROX$(\hat\kappa)$ returns the origin at every stage $j\le33$ for
$\hat\kappa\in\{1,2\}$. For
$\hat\kappa=4$, the points $(0,0)$ and $(1,1)$ both minimize $Q_\eta$ at
stages $0$ and $1$, and PROX returns the origin at every stage $j\le33$ when
ties are broken towards the current center. These stages cover the budgets
of DISC: the least common denominator of Section~\ref{sec:exact-data} is
$\Delta=128$, so $R=2^{23}$, $\tau=2^{-26}$ and $J_{\hat\kappa}=28,28,29$ for
$\hat\kappa=1,2,4$. Lemma~\ref{lem:diagcert} rejects the origin:
$H+\Lambda=\bigl(\begin{smallmatrix}191/64&-3\\-3&191/64\end{smallmatrix}\bigr)$
has eigenvalue $-\frac1{64}$. At $(1,1)$ the matrix has diagonal
$\frac{193}{64}$ and is positive definite, so $F\in\mathfrak D$, and DISC
stops at a later guess by Theorem~\ref{thm:diagdiscovery}(b).
\end{remark}
